\documentclass[11pt]{Article}
% \usepackage{fullpage}
\usepackage{graphicx}
\usepackage{amsmath}
\usepackage{pdfpages}
\usepackage{lipsum}
\usepackage{blindtext} 
\usepackage{fancyhdr}
\usepackage[colorinlistoftodos]{todonotes}


%syntax highlighting and code blocks
\usepackage{listings}
% \usepackage{minted}

\usepackage{tikz}
\usetikzlibrary{positioning,calc}

% Language setting
% Replace `english' with e.g. `spanish' to change the document language
% \usepackage[english]{babel}
\usepackage{multicol}
\usepackage{geometry}
\usepackage{caption}
\usepackage{float}
\usepackage[hidelinks]{hyperref}
\usepackage[backend=biber,style=apa,autocite=inline]{biblatex}
\DeclareLanguageMapping{english}{english-apa}
\setlength{\columnsep}{1cm}
\geometry{
        a4paper,
        total={170mm,257mm},
        left=20mm,
        top=20mm,
        }

\usepackage{listings}
\usepackage{xcolor}

%%%%%%%%%%%%%%%%%%%%%%%%%%%%%%%%%%%%%%%%%%%%%%%%%%%%%%%%%%%%%%%%%%%%%%%%%%%%%%%%%%%%%%%%%%%%%%%%
%%%%%%%%%%%%%%%%%%%%%%%%%%%%%%%%%%%%%%%%%%%%%%%%%%%%%%%%%%%%%%%%%%%%%%%%%%%%%%%%%%%%%%%%%%%%%%%%


\definecolor{codegreen}{rgb}{0,0.6,0}
\definecolor{codegray}{rgb}{0.5,0.5,0.5}
\definecolor{codepurple}{rgb}{0.58,0,0.82}
\definecolor{backcolour}{rgb}{0.95,0.95,0.92}

\lstdefinestyle{mystyle}{
    backgroundcolor=\color{backcolour},   
    commentstyle=\color{codegreen},
    keywordstyle=\color{blue},
    numberstyle=\tiny\color{codegray},
    stringstyle=\color{codepurple},
    basicstyle=\ttfamily\footnotesize,
    breakatwhitespace=false,         
    breaklines=true,                 
    captionpos=b,                    
    keepspaces=true,                 
    numbers=left,                    
    numbersep=5pt,                  
    showspaces=false,                
    showstringspaces=false,
    showtabs=false,                  
    tabsize=2
} 

%%%%%%%%%%%%%%%%%%%%%%%%%%%%%%%%%%%%%%%%%%%%%%%%%%%%%%%%%%%%%%%%%%%%%%%%%%%%%%%%%%%%%%%%%%%%%%%%
%%%%%%%%%%%%%%%%%%%%%%%%%%%%%%%%%%%%%%%%%%%%%%%%%%%%%%%%%%%%%%%%%%%%%%%%%%%%%%%%%%%%%%%%%%%%%%%%


% Useful packages
% \usepackage{amsmath}
% \usepackage{siunitx}
% \usepackage{wrapfig}
% \usepackage{float}
% \usepackage{graphicx}
% \usepackage{subcaption}
% \usepackage[colorlinks=true, allcolors=blue]{hyperref}
% \usepackage{xcolor}

\graphicspath{
    {img/}
}
\bibliography{./bib.bib}

\newenvironment{Figure}
  {\par\medskip\noindent\minipage{\linewidth}}
  {\endminipage\par\medskip}
\newcommand\C[1]\null

%remove if done, this extends the page to see comments
% \paperwidth=\dimexpr \paperwidth + 8cm\relax
% \oddsidemargin=\dimexpr\oddsidemargin + 3cm\relax
% \evensidemargin=\dimexpr\evensidemargin + 3cm\relax
% \marginparwidth=\dimexpr \marginparwidth + 3cm\relax


\title{Hac-man: Game design document}

\author{%
  \begin{tabular}{c c c}
    Mirko van de Hoef   &    Tom Selier\\
    \texttt{m.a.j.vandeHoef@students.uu.nl} & \texttt{t.b.j.selier@students.uu.nl} \\
    Student Number: 2176777 & Student Number: 5412498 
  \end{tabular}
}


\usepackage{eso-pic}
\usepackage{atbegshi}
\AtBeginShipoutFirst{\AddToShipoutPictureBG*{
    \begin{tikzpicture}[overlay, remember picture, outer sep=0pt, inner sep=1ex]
      \node [anchor=north west] at ($(current page.north)-(.5*\textwidth-1ex, \headheight)$)
      {\parbox{10cm}{UU University Utrecht\\ Functional Programming\\2024, Utrecht, Netherlands}};
    \end{tikzpicture}
}}




%%%%%%%%%%%%%%%%%%%%%%%%%%%%%%%%%%%%%%%%%%%%%%%%%%%%%%%%%%%%%%%%%%%%%%%%%%%%%%%%%%%%%%%%%%%%%%%%
%%%%%%%%%%%%%%%%%%%%%%%%%%%%%%%%%%%%%%%%%%%%%%%%%%%%%%%%%%%%%%%%%%%%%%%%%%%%%%%%%%%%%%%%%%%%%%%%
%beginning of document
\begin{document}
\lstset{style=mystyle}
\maketitle
% \begin{abstract}
% \lipsum[1]
% \end{abstract}
% {\bf \textit{Keywords--}} game game gamin
%%%%%%%%%%%%%%%%%%%%%%%%%%%%%%%%%%%%%%%%%%%%%%%%%%%%%%%%%%%%%%%%%%%%%%%%%%%%%%%%%%%%%%%%%%%%%%%%
%%%%%%%%%%%%%%%%%%%%%%%%%%%%%%%%%%%%%%%%%%%%%%%%%%%%%%%%%%%%%%%%%%%%%%%%%%%%%%%%%%%%%%%%%%%%%%%%

%introduction
\section{Introduction} \label{ch:Introduction}
The objective of this design document is to present the design choices for the creation of Pac-Man, along with their respective justifications. The decision to implement Pac-man is based on on the expected freedom the game offers in its implementation. Furthermore, Pac-Man will allow us to implement some of the optional requirements in addition to the minimum requirements. Additionally, some variations on the fundemental gameplay are introduced to make the game more unique.

\section{Variations on the base game} \label{ch:twist}
In order to introduce a new learning opportunity and make the game more unique, a player adversary will be introduced. This will introduce, a co-op type, multiplayer aspect to the game. This feature will introduce a second set of player inputs which allows the control of one of the ghosts.

\section{Data types} \label{ch:dataTypes}
%ghost and player

The first thing to implement are models for the player and the ghosts. These objects need to move, and store some information about them. 

Movable objects in the game are called entities. These entities have a direction in which they are moving, a speed and a current position. These three data types are wrapped together in a record called movement. Furthermore, an entity can be alive or dead. 

Then, we can create a data type for the player and the ghosts. A player gets lives, next to being an entity. And the ghosts get some identifiers for it's name and it's current behaviour. Lastly, a data type got created for an optional local MultiPlayer feature, where a second player can control one of the ghosts.

\begin{lstlisting}[language=Haskell]
  data isAlive = Alive | Dead

  data Direction = Left | Right | Up | Down
  
  type EntityPosition = (Float, Float)
  
  data Movement {
    direction :: Direction,
    speed :: Int,
    position :: EntityPosition
  }

  data Entity {
    movement :: Movement,
    alive :: isAlive
  }

  data GhostType = Inky 
                 | Pinky 
                 | Blinky 
                 | Clyde 
  
  data BehaviourMode = Chase | Scatter | Frightened | Home
  type Lives = Int

  data Player = {
    entity :: Entity,
    lives :: Lives
  }

  data Ghost = {
    entity :: Entity,
    ghostName :: GhostType,
    behaviourMode :: Behaviour
  }

  data PlayerGhost = {
    entity :: Entity,
    ghostName :: GhostType
  }
\end{lstlisting}

% levels
\noindent The maze consists of two type of tiles, where each tile is an x by x square in a grid. The tile can be either a floor or a wall where each floor can be either one of the consumbles or an empty floor. The walls can be either a corner or a straight wall, each with their own orientation. This will be implemented as follows:
\begin{lstlisting}[language=Haskell]
  type Level = Int

  data CornerOrientation = SE | NW | NE | SW
  data WallOrientation   = Horizontal | Vertical
  
  data WallShape  = Corner CornerOrientation
                  | Wall WallOrientation 

  data ConsumableType = Pellet | SuperPellet | Cherry

  data FloorType = Consumable ConsumableType | EmptyTile
  data Tile = Floor FloorType | Wall WallShape
\end{lstlisting}


The maze will utilise a Map structure in order to speed up execution. This would result in a reduction in the scaling from $O(n)$ to $O(logn)$ for tile lookup. As the program is expected to lookup a tile every time one has to check for collision, it is crucial that this takes as little time as possible. This is only possible if the key, which is the x and y coordinates of the tile, can be ordered. The instance below is an example but is important to show that the keys are indeed orderable. In this case it is possible using the implementation below. Thus, two tiles in the maze might look like: [((1, 0), Floor (Consumable Pellet)), ((2, 0), Wall (Corner SE))]. 

\begin{lstlisting}[language=Haskell]
  type TilePosition = (Int, Int) -- (x, y)

  type Maze = Map TilePosition Tile

  instance Ord TilePosition where 
    (<=) (x, y) (a, b) | x < a = True 
                       | x > a = False
                       | Otherwise = y <= b
\end{lstlisting}

\newpage
%gamestate
The game state is the final data type that contains all the data of the game. It's job is to be the model in the MVC structure. 

Some data types that have been added here, but not discussed before are some scorekeepers, like elapsedTime, score and level. Furthermore, the current status of the game, e.g., it being paused or running, is kept in the game state.
\begin{lstlisting}[language=Haskell]
  type Name = String

  type Sprite = Gloss.Picture
  
  type Sprites = Map.Map Name Sprite
  
  data GameStatus = Running | GameOver | Paused | Quitting deriving (Eq, Show)
  
  newtype WindowInfo = MkWindowInfo {resolution :: (Int, Int)}
  
  type HighScores = [String]
  
  data GameState = MkGameState
    { status      :: GameStatus
    , maze        :: Maze
    , isNewMaze   :: Bool --boolean flag if the Maze has changed or not
    , oldMaze     :: Gloss.Picture
    , sprites     :: Sprites --map of sprites
    , elapsedTime :: Float
    , deltaTime   :: Float
    , enableDebug :: Bool
    , windowInfo  :: WindowInfo --resolution of window
    , player      :: Player
    , ghosts      :: [Ghost]
    , pelletC     :: (Int, Int) -- total vs eaten
    , highscores  :: HighScores -- list of highscores (retrieved from file)
    , level       :: Int        -- current level
    }
  
\end{lstlisting}


\section{Minumum Requirements} \label{ch:minumumRequirements}
\noindent \textbf{Player} The player can control Pac-Man, and the menus, using the WASD-keys on the keyboard.\\

\noindent \textbf{Enemies} The computer controls the red ghost, called Blinky. Blinky is the simplest of four ghost and tries to chase Pac-Man.\\

\noindent \textbf{Randomness} A ghost can exist in three states; Scatter, chase and frightened. When a ghost is frightened it will pick a random direction at each intersection.\\

\noindent \textbf{Animation} The amount of animations in Pac-Man is limited. The animations that will be implemented are:
\begin{itemize}
  \item Pac-Man's walking animation
  \item Pac-Man's death animation
  \item Ghost's movement
  \item Ghost's frightened\\
\end{itemize}

\noindent \textbf{Pause} The game can be paused using the "pause" or "esc" keys on the keyboard.\\

\noindent \textbf{Interaction with the file system} A running highscore will be kept in a file.\\


\newpage
\section{Optional Requirements} \label{ch:optionalRequirements}

\textbf{Levels} Originally, Pac-Man only contains one level, that gets more difficult as you complete them. At least one additional level will be added as an option for the player. In the chapter \ref{ch:dataTypes}, it was discussed how the maze data type is made. This type lends itself for easy creation of levels as one needs to simply provide each coordinate (written in tuple), with the desired tile type. Another option would be to write the levels as a single list of special characters, and making a function that converts to the key and value for the Map. The storage of each level may be accomplished through the utilisation of either the file system that parses the levels or the implementation of hard-coded levels.

\textbf{Different enemies} Pac-Man contains several ghosts. At a minimum Blinky will be implemented. Additional ghosts will be Pinky (pink), Inky (cyan) and Clyde (yellow), and if time allows, some additional ghosts could be added.

\textbf{Multi-player} The original way to play against another player in Pac-Man is to beat eachother's highscore. Optionally, a way to control one of the ghosts will be added to the game. In order to simplify the process, the entities and player controller will be defined in the most abstract manner possible, thus facilitating the reuse of fucntions. For example, a "moveStep" function which could take in a current position and a numerical to change the position of an entity. As seen in chapter \ref{ch:dataTypes}, the ghosts and player use the same movement data type, thus facilitating the use of these generelized functions.

\begin{lstlisting}
  moveStep :: Entity -> Int -> Entity
\end{lstlisting}


\section{Pure/Impure seperation} \label{ch:pureSeperation}
The goal is to go from impure to pure parts as soon as possible. To do this, we need to create functions that take only inputs from its parameters. For example, the highscore is kept in a file. Reading/writing this file uses impure function, but updating the list of highscores will be pure. Thus, reading/ writing will be impure but updating the highscore list is pure.

\begin{lstlisting}
  saveHighscore ::  [Int] -> IO ()
  updateHighscore :: [Int] -> Int -> [Int]
\end{lstlisting}

\section{Abstraction} \label{ch:abstraction}
Several steps for abstraction have already been taken. For example, the "moveStep" function takes an entity, and does not care which entity (player or NPC) it is actually a part of. Besides movement, entities are also alive or dead, using the data type "isAlive" 

Furthermore, since tiles are both consumables or empty, it should be trivial to implement a function that updates a tile based on the position of the player. An example function might be

\begin{lstlisting}
  updateTile :: Player -> Tile -> Tile
\end{lstlisting}

\end{document}


% \begin{lstlisting}[language=Haskell]
%   sum :: [a] -> a
%   sum = foldr (+) 0
% \end{lstlisting}

% \begin{Figure}
%   \centering
%   \includegraphics[scale=0.45]{placeholder.png}
%   \captionof{figure}{An example}
%   \label{fig:placeholder}
% \end{Figure}


% \begin{Figure}
%   \centering
%   \includegraphics[scale=0.45]{Bubble sensor diagram.png}
%   \captionof{figure}{Schematical overview of the hemispherical pressure sensor. [I] Silicone layer, [II] barometric sensor, [III] Working fluid, [IV] Adhesive, [V] Flex-PCB.}
%   \label{fig:SchematicalOverviewHemisphere}
% \end{Figure}
